\documentclass[UTF8,fontset=windows,12pt]{ctexart}
%\setcounter{tocdepth}{1}
%\setcounter{secnumdepth}{4}
%\usepackage{ctex} %若要输入中文请取消注释或者把article改成ctexart
% This first part of the file is called the PREAMBLE. It includes
% customizations and command definitions. The preamble is everything
% between \documentclass and \begin{document}.
%\usepackage[OT1]{fontenc}
%\usepackage{fontspec}
%\setmainfont{Times New Roman}
%\setCJKmainfont{Microsoft YaHei}
\usepackage[left=1.91cm,right=1.91cm,top=2.54cm,bottom=2.54cm]{geometry}
\usepackage{amsfonts}
\usepackage{amsmath}
\usepackage{amssymb}
\usepackage[upint,lcgreekalpha]{stix}
\usepackage[type1]{garamondlibre}
%\usepackage{esint}
\usepackage{pgfplots}
\pgfplotsset{compat=1.18}
\usepackage{tikz}
\usetikzlibrary{arrows}
\usepackage{extarrows}
\usepackage[only,llbracket,rrbracket]{stmaryrd}
\usepackage{titlesec}
\usepackage{bm}
\usepackage{graphicx}
\usepackage{appendix}
\usepackage{tikz-cd}
\usepackage{tikz-3dplot}
\usetikzlibrary{arrows.meta, decorations.pathreplacing, calc, patterns}
\usepackage{float}
\usepackage{mathtools}
\usepackage{amsthm}
\usepackage{verbatim}              % better theorem environments
\usepackage{setspace}
\usepackage{enumitem}
\setenumerate[1]{itemsep=0pt,partopsep=0pt,parsep=\parskip,topsep=5pt}
\setitemize[1]{itemsep=0pt,partopsep=0pt,parsep=\parskip,topsep=5pt}
\setdescription{itemsep=0pt,partopsep=0pt,parsep=\parskip,topsep=5pt}
\usepackage{xcolor}
\usepackage[colorlinks,linkcolor=black,anchorcolor=black,citecolor=black]{hyperref}
%\usepackage{apacite}
\usepackage[numbers]{natbib}
%\usepackage{showkeys}
\usepackage{cases}
\usepackage{fancyhdr}
\usepackage[scaled=0.92]{helvet}	% set Helvetica as the sans-serif font
%\usepackage{upgreek}
%\renewcommand{\rmdefault}{ptm}		% set Times as the default text font
%\usepackage{newtxtext}
\bibliographystyle{plain}
% If AMS-LaTeX is used, it can be loaded before or after mtpro2		
% The following loads metro and defines some common MTPro options [2, 4]
%\usepackage[lite,subscriptcorrection,nofontinfo]{mtpro2}
%\pdfmapfile{+mtpro2.map}

%\usepackage{txfonts}

% various theorems, numbered by section

\makeatletter
\renewenvironment{proof}[1][\proofname]{\par
  \pushQED{\qed}%
  \normalfont \topsep6\p@\@plus6\p@\relax
  \trivlist
  \item[\hskip\labelsep
        \bfseries % 关键修改：斜体改为加粗
        #1\@addpunct{.}]\ignorespaces
}{%
  \popQED\endtrivlist\@endpefalse
}
\makeatother
\newtheoremstyle{kaiti-theorem}   % 样式名称
  {3pt}                           % 上方间距
  {3pt}                           % 下方间距
  {\kaishu\upshape}               % 正文字体：中文楷体 + 英文直立
  {}                              % 缩进量
  {\bfseries\upshape}      % 头部字体（定理名）：中文楷体粗体 + 英文直立粗体
  {.}                             % 定理名后标点
  {0.5em}                         % 定理名后间距
  {}                              % 头部说明（留空）

% 应用自定义样式
\theoremstyle{kaiti-theorem}
\newtheorem{thm}{定理}[section]
\newtheorem{nota}[thm]{记号}
\newtheorem{question}{问题}
\newtheorem*{questionn}{思考}
\newtheorem{exe}{习题}[section]
\newtheorem{assump}[thm]{假设}
\newtheorem*{thmm}{定理}
\newtheorem{defn}{定义}[section]
\newtheorem{lem}[thm]{引理}
\newtheorem*{lemm}{引理}
\newtheorem{prop}[thm]{命题}
\newtheorem*{propp}{命题}
\newtheorem*{claim}{断言}
\newtheorem{cor}[thm]{推论}
\newtheorem{conj}[thm]{猜想}
\newtheorem{rmk}{注记}[section]
\newtheorem{ex}{例}[section]
\newtheorem{pf}{证明}
\newtheorem{problem}{作业题}
\numberwithin{equation}{section}





\DeclareMathOperator{\id}{id}
\DeclareMathOperator*{\esssup}{ess\,sup}
\renewcommand{\Re}{\textbf{Re}\,}  
\renewcommand{\Im}{\textbf{Im}\,}  
\newcommand{\R}{\mathbb{R}}  
\newcommand{\FH}{\mathbb{H}}  
\newcommand{\C}{\mathbb{C}}  
\newcommand{\Z}{\mathbb{Z}}
\newcommand{\X}{\mathbb{X}}
\newcommand{\N}{\mathbb{N}}
\newcommand{\Q}{\mathbb{Q}}
\newcommand{\T}{\mathbb{T}}
\newcommand{\TT}{\mathcal{T}}
\newcommand{\TP}{\overline{\partial}{}}
\newcommand{\TJ}{\langle\TP\rangle}
\newcommand{\TL}{\overline{\Delta}{}}
\newcommand{\curl}{\text{curl }}
\newcommand{\dive}{\text{div }}
\newcommand{\bd}[1]{\mathbf{#1}}  % for bolding symbols
\newcommand{\bds}[1]{\boldsymbol{#1}}  % for bolding symbols
\newcommand{\RR}{\mathcal{R}}      % for Real numbers
\newcommand{\sh}{\mathcal{S}}  
\newcommand{\sph}{\mathbb{S}}  
\newcommand{\Om}{\Omega}
\newcommand{\OmT}{{\Omega_T}}
\newcommand{\ut}{{U_T}}
\newcommand{\vp}{\varphi}
\newcommand{\Th}{\Theta}
\newcommand{\Lam}{\Lambda}
\newcommand{\Omb}{{\overline{\Omega}}}
\newcommand{\OmbT}{{\overline{\Omega_T}}}
\newcommand{\ub}{{\overline{U}}}
\newcommand{\ubt}{{\overline{U_T}}}
\newcommand{\q}{\quad}
\newcommand{\p}{\partial}
\newcommand{\pb}{\boldsymbol{\partial}}
\newcommand{\lap}{\Delta}
\newcommand{\dd}{\mathfrak{D}}
\newcommand{\DD}{\mathcal{D}}
\newcommand{\den}{{\bar\rho}_{0}}
\newcommand{\Dt}{{\p}_{t}+v^{k}{\p}_{k}}
\newcommand{\col}[1]{\left[\begin{matrix} #1 \end{matrix} \right]}
\newcommand{\noi}{\noindent}
\newcommand{\nab}{\nabla}
\newcommand{\cnab}{\overline{\nab}}
\newcommand{\cp}{\overline{\partial}{}}
\newcommand{\dx}{\,\mathrm{d}x}
\newcommand{\dxx}{\,\mathrm{d}\bds{x}}
\newcommand{\xxi}{\boldsymbol{\xi}}
\newcommand{\eeta}{\boldsymbol{\eta}}
\newcommand{\dxxi}{\,\mathrm{d}\boldsymbol{\xi}}
\newcommand{\deeta}{\,\mathrm{d}\boldsymbol{\eta}}
\newcommand{\dxi}{\,\mathrm{d}\xi}
\newcommand{\deta}{\,\mathrm{d}\eta}
\newcommand{\dy}{\,\mathrm{d}y}
\newcommand{\dyy}{\,\mathrm{d}\bds{y}}
\newcommand{\dz}{\,\mathrm{d}z}
\newcommand{\dzz}{\,\mathrm{d}\bds{z}}
\newcommand{\dph}{\,\mathrm{d}\varphi}
\newcommand{\dt}{\,\mathrm{d}t}
\newcommand{\dr}{\,\mathrm{d}r}
\newcommand{\dtt}{\,\mathrm{d}\tau}
\newcommand{\dS}{\,\mathrm{d}S}
\newcommand{\dss}{\,\mathrm{d}\sigma}
\newcommand{\dmu}{\,\mathrm{d}\mu}
\newcommand{\dnu}{\,\mathrm{d}\nu}
\newcommand{\dV}{\,\mathrm{d}V}
\newcommand{\ds}{\,\mathrm{d}s}
\newcommand{\drr}{\,\mathrm{d}\rho}
\newcommand{\dist}{\text{dist }}
\newcommand{\vol}{\text{vol}\,}
\newcommand{\volo}{\text{vol}(\Omega)}
\newcommand{\spt}{\text{Spt}\,}
\newcommand{\Tr}{\text{Tr}\,}
\newcommand{\lee}{\langle}
\newcommand{\ree}{\rangle}
\newcommand{\xee}{\langle\xi\rangle}
\newcommand{\xxee}{\langle\boldsymbol{\xi}\rangle}
\newcommand{\xeta}{\langle\eta\rangle}
\newcommand{\xeeta}{\langle\boldsymbol{\eta}\rangle}
\newcommand{\kk}{\kappa}
\newcommand{\lam}{\lambda}
\newcommand{\io}{\int_{\Omega}}
\newcommand{\iu}{\int_{U}}
\newcommand{\is}{\int_{\Sigma}}
\newcommand{\ipo}{\int_{\partial\Omega}}
\newcommand{\ipu}{\int_{\partial U}}
\newcommand{\ig}{\int_{\Gamma}}
\newcommand{\ir}{\int_{\R}}
\newcommand{\ird}{\int_{\R^d}}
\newcommand{\irn}{\int_{\R^n}}
\newcommand{\fpi}{\frac{1}{(2\pi)^{\frac{d}{2}}}}
\newcommand{\ddt}{\frac{\mathrm{d}}{\mathrm{d}t}}
\newcommand{\ddx}{\frac{\mathrm{d}}{\mathrm{d}x}}
\newcommand{\eps}{\varepsilon}
\providecommand{\jump}[1]{\left\llbracket #1 \right\rrbracket }
\providecommand{\len}[1]{\lee #1 \ree }
\providecommand{\ino}[1]{\left\| #1 \right\| }
\providecommand{\bno}[1]{\left| #1 \right| }
\newcommand{\red}{\textcolor{red}}
\newcommand{\green}{\textcolor{green}}
\newcommand{\blue}{\textcolor{blue}}
\newcommand{\nnr}{{[n+1]}}
\newcommand{\nnn}{{[n]}}
\newcommand{\nnl}{{[n-1]}}
\newcommand{\nnll}{{[n-2]}}	
\newcommand{\mmr}{{[m+1]}}
\newcommand{\mmm}{{[m]}}
\newcommand{\mml}{{[m-1]}}
\newcommand{\mmll}{{[m-2]}}

%---------------Special variables--------------
\newcommand{\CC}{\mathfrak{C}}  
\newcommand{\NN}{\mathbf{N}}
\newcommand{\LH}{\mathcal{L}}
\newcommand{\HH}{\mathcal{H}}  
\newcommand{\EE}{\mathcal{E}}  
\newcommand{\FT}{\mathcal{F}}  
\newcommand{\fT}{\mathbf{T}}  
\newcommand{\FA}{\mathbf{A}}  
\newcommand{\MH}{\mathcal{M}}
\newcommand{\NH}{\mathcal{N}}
\newcommand{\PH}{\mathcal{P}}
\newcommand{\VH}{\mathcal{V}}
\newcommand{\XX}{\mathfrak{X}}
\newcommand{\xx}{{\bds{x}}}
\newcommand{\bx}{\mathbf{x}}
\newcommand{\by}{\mathbf{y}}
\newcommand{\bp}{\mathbf{p}}
\newcommand{\bq}{\mathbf{q}}
\newcommand{\pp}{{\bds{p}}}
\newcommand{\qq}{{\bds{q}}}
\newcommand{\uu}{\mathbf{u}}
\newcommand{\ww}{\mathbf{w}}
\newcommand{\vv}{\mathbf{v}}
\newcommand{\yy}{{\bds{y}}}
\newcommand{\zz}{{\bds{z}}}
\newcommand{\zo}{\mathbf{0}}
\newcommand{\ee}{\mathbf{e}}
\newcommand{\fa}{\mathbf{a}}
\newcommand{\fb}{\mathbf{b}}
\newcommand{\fc}{\mathbf{c}}
\newcommand{\ff}{\bds{f}}
\newcommand{\fr}{\mathbf{r}}
\newcommand{\fh}{\mathbf{h}}
\newcommand{\fm}{\mathbf{m}}
\newcommand{\fg}{\mathbf{g}}
\newcommand{\FF}{\mathbf{F}}
\newcommand{\GG}{\mathbf{G}}
\newcommand{\BB}{\mathbf{B}}



\begin{document}
\title{\LARGE{\bf  2026年春季学期现代偏微分方程作业三}}
\author{Pohozaev恒等式、线性抛物方程、消失粘性法}
\date{截止时间：2026.5.7下课前（电子版截至当天23:59:59）}

\maketitle


\boxed{\textbf{以下两题均可忽略光滑逼近的步骤，且可以默认第一题无穷远处的边界积分皆为零。}}

\begin{problem}[习题2.8.1]
$\R^d$中的质量临界Schr\"odinger方程的基态 (ground state) $Q$是满足如下半线性椭圆方程的\underline{径向正解}，且在无穷远处衰减到零
\[
-\lap Q + Q - |Q|^{\frac{4}{d}}Q =0 ,\q\q \xx\in\R^d.
\]
\begin{itemize}
	\item [(1)] 模仿定理2.8.2的方法证明: $\|Q\|_{L^{2+\frac{4}{d}}(\R^d)}^{2+\frac{4}{d}}=(1+\frac{d}{2})\|Q\|_{L^2(\R^d)}^2$.
	\item [(2)] 已知如下 Galiargo-Nirenberg 插值不等式成立，且$u=Q$时该不等式取等号
	\[
	\|u\|_{L^{2+(4/d)}(\R^d)}^{2+(4/d)}\leqslant C\|u\|_{L^2(\R^d)}^{4/d}\|\nabla u\|_{L^2(\R^d)}^2.
	\]写出此时的最佳常数$C$关于$\|Q\|_{L^2(\R^d)}^2$的表达式。
\end{itemize}
\end{problem}

\begin{problem}[习题2.8.2]  设有界区域$U\subset\R^d$边界光滑，$V\in C^1(\ub)$是给定的位势函数，$\lam>0$是$(-\lap+V(\xx)I)$算子(带Dirichlet零边值)的一个特征值，$u\in H^2(U)\cap H_0^1(U)$是对应的特征函数，即
\[
-\lap u +V(\xx) u= \lam u\text{ in }U,\q\q u|_{\p U}=0.
\]证明：
\[
\lam\iu u^2\dxx=\frac12\ipu (\xx\cdot N(\xx))\left(\frac{\p u}{\p N}\right)^2\dS_\xx +\frac12 \iu (2V(\xx)+\xx\cdot\nab V(\xx))u^2\dxx.
\]
\end{problem}



%\begin{problem}[习题2.2.4]
%考虑具有 Robin 边界条件的 Poisson 方程
%\[
%-\lap u=f\text{ in }U,\q\q\q u+\frac{\p u}{\p N}=0\text{ on }\p U.
%\]请定义该问题的弱解 $u\in H^1(U)$，并讨论对于给定 $f\in L^2(U)$ 时解的存在性和唯一性。
%
%提示：证明双线性型的强制性可以使用类似于定理 1.4.5 中 Poincar\'e 不等式证明的技巧。
%\end{problem}

\begin{problem}[习题3.2.1]
设 $f\in L^2(U)$ ，$u_m=\sum\limits_{k=1}^md_m^k w_k$ 满足 
\[
\int_U\nab u_m\cdot\nab w_k\dxx=\int_U f\, w_k\dxx,\quad \forall 1\leqslant k\leqslant m,
\] 其中 $\{w_k\}$ 是 $H_0^1(U)$ 的标准正交基。证明：$\{u_m\}$ 存在子列在 $H_0^1(U)$ 中弱收敛于Poisson方程 $-\lap u=f ~(\xx\in U)$, $u|_{\p U}=0$ 的弱解 $u$.
\end{problem}


\begin{problem}[习题3.2.2+问题3.4.1]
设 $g\in L^2(U)$ 并设 $u$ 是如下问题的光滑解：
\begin{equation}
\p_t u - \lap u =0~\text{ in }U_T,\quad u=0~\text{ on }[0,T]\times\p U,\quad u=g~\text{ on }\{t=0\}\times U.
\end{equation}
\begin{itemize}
\item [(1)] 证明：对任意的$t\geqslant 0$有$\|u(t,\cdot)\|_{L^2(U)}\leqslant e^{-\lam_1 t}\|g\|_{L^2(U)}.$ 这里 $\lam_1>0$ 是具有零 Dirichlet 边界条件的 $-\lap$ 在 $U$ 上的主特征值。
\item [(2)] 如果再假设 $g\in C^1(\ub)$ 且满足相容性条件 $g|_{\p U}=0$，证明：存在常数 $C>0$ 使得对任意 $t>0$ 都有 $\sup\limits_{\xx\in \ub}|u(t,\xx)|\leqslant Ce^{-\lam_1 t}.$
\end{itemize}
提示：对(2)，如果初值是对应于$(-\lap)$算子(带Dirichlet边界条件时)的主特征值 $\lam_1$ 的特征函数 $w_1$，那对应热方程的解是什么？然后用问题2.6.2的结论。
\end{problem}

%\begin{problem}[习题3.4.1]
%设 $u$ 为热方程 $\p_t u-\lap u+cu=0$ 在 $\R_+\times U$ 中的光滑解，在 $\{t=0\}\times \ub$ 上 $u=g$，边界条件为在 $[0, \infty) \times \p U$ 上 $u=0$。这里 $g(\xx),c(t,\xx)$ 是给定的连续函数。
%\begin{itemize}
%\item [(1)] 若存在常数 $\gamma$ 使得 $c\geqslant \gamma>0$，证明：存在常数 $A>0$ 使得对任意的 $(t,\xx)\in(0,T]\times U$，均有 $|u(t,\xx)|\leqslant A e^{-\gamma t}$. 
%\item [(2)] 当 $g\geqslant 0$ 且 $c$ 有界时，证明 $u$ 是非负的。
%\end{itemize}
%\end{problem}

\begin{problem}[问题3.5.1]
记 $\Om = \R_+\times\R\subset\R^2$，其边界为 $\p\Om = \{ x_1 = 0 \}$. 设向量场 $\mathbf{b}=(b_1,b_2) \in W^{3,\infty}([0,T]\times\Omb;\R^2)$. 任给参数 $\eps \in(0,1]$，考虑粘性对流--扩散方程初边值问题：
\begin{equation}\label{CDE}\tag{$P_\eps$}
\begin{cases}
\partial_t u^\eps + \mathbf{b}(t,\boldsymbol{x}) \cdot \nabla u^\eps - \eps \Delta u^\eps = 0&t\in(0,T),\xx\in\Om, \\
\p_1 u^\eps|_{\p\Om}=0& t\in[0,T],\xx\in\p\Om,\\
u^\eps(0,\boldsymbol{x}) = g(\boldsymbol{x}) \in H^2(\Om)&\xx\in\Omb.
\end{cases}
\end{equation}
已知对任意 $\eps > 0$，\eqref{CDE}存在充分光滑的解 $u^\eps$. 
\begin{itemize}
\item [(1)] 若 $b_1=0$. 证明如下结论。
	\begin{itemize}
	\item [(1a)] $\sup\limits_{t\in[0,T]}\|u^\eps\|_{H^1(\Om)}$ 关于 $\eps$ 一致有界，进而 $\{u^\eps\}$ 存在子列在 $L^\infty(0,T;H^1(\Om))$ 中弱-*收敛到 $u\in L^\infty(0,T;H^1(\Om))$, 且 $u$ 是 无粘传输方程 $\partial_t u + \mathbf{b}(t,\boldsymbol{x}) \cdot \nabla u =0$ 的弱解，初值仍为 $g$. 
	\item [(1b)] 此时 $\p_1 u|_{\p\Om}$ 是否一定为零？
	\item [(1c)] $\sup\limits_{t\in[0,T]}\|u^\eps-u\|_{L^2(\Om)}\leqslant C\sqrt{\eps}$. 
	\end{itemize}
\item [(2)] 若 $b_1=1$ 且 $b_2$ 不依赖 $x_1$. 证明：$\sup\limits_{t\in[0,T]}\|u^\eps\|_{L^2(\Om)}$ 关于 $\eps$ 仍然一致有界，并问：$\sup\limits_{t\in[0,T]}\|u^\eps\|_{H^1(\Om)}$ 关于 $\eps$ 是否一致有界？证明结论或者构造反例。
%\item [(3)] 如果把初值条件降为$g\in H^1(\Om)$，(2)的其它条件不变，则(2)的结论如何？
\end{itemize}
提示：若$g\in H^2(\Om)$且方程解充分光滑，则$g$应该满足相容性条件$\p_1 g|_{\p\Om}=0$，但在$g\in H^1(\Om)$时无法定义$\p g$的迹。对(2)，思考如何对$u^\eps$的一阶导数做一致估计。
\end{problem}

\begin{problem}[补充题]
	设维数$d\geqslant 2$，给定$f\in L^2(U)$. 对常数$\eps>0$，考虑如下带有钳支边界条件的奇异摄动四阶椭圆方程
	\[
	\eps\lap^2u_\eps - \lap u_\eps = f \text{ in }U,\quad\quad u_\eps=\frac{\p u_\eps}{\p N}=0\text{ on }\p U.
	\]已知对任意给定的 $\eps > 0$，该方程存在唯一的弱解 $u_\eps\in H_0^2(U)$.
	       \begin{itemize}
   	\item [(1)] 证明：存在 $u_0 \in H_0^1(U)$，使得当 $\eps \to 0^+$ 时，$u_\eps\xrightarrow{H_0^1(U)} u_0$，并求出 $u_0$ (在弱意义下)满足的极限方程。
   	\item [(2)] 证明：$\lim\limits_{\eps \to 0^+} \eps \int_U |\lap u_\eps|^2 \dxx = 0$.
   	\item [(3)] 若$f$使得(1)中所得的$u_0\in H^2(U)$，且 $\frac{\p u_0}{\p N}$ 在 $\partial U$ 上不恒为零。证明：$\lim\limits_{\eps\to 0_+}\|u_\eps\|_{H^2(U)}=+\infty$.
   	\item [(4)] (选做) 仍然假设$f$使得(1)中所得的$u_0\in H^2(U)$，而且极限函数还恰好满足 $\frac{\p u_0}{\p N}$在 $\partial U$ 上恒为零，回答如下问题：
   	\begin{itemize}
   		\item [(4a)] 证明：当$\eps\to 0_+$时，$u_\eps\xrightarrow{H^2}u_0$，且存在不依赖$\eps$的常数$C>0$使得 $\|u_\eps - u_0\|_{H_0^1(U)} \leqslant C\sqrt{\eps}$. 
   		\item [(4b)] 若进一步假设$f\in H_0^1(U)$，证明： 存在不依赖$\eps$的常数$C''>0$使得 $\|u_\eps - u_0\|_{H^2(U)} \leqslant C''\sqrt{\eps}$. 
   		\item [(4c)] 如果去掉$f\in H_0^1(U)$，此时是否还存在$a>0$，使得存在不依赖$\eps$的常数$C'>0$满足$\|u_\eps - u_0\|_{H^2(U)} \leqslant C'\eps^a$？证明你的结论。
   	\end{itemize}
   \end{itemize}
   
   提示：(3) 联想迹定理的结论；(4) 注意对 $f\in H_0^2$,  $\|\lap f\|_{L^2}$ 和 $\|f\|_{H^2}$ 是等价范数；对 (4c) 可以考虑模仿问题 2.2.2 构造高频序列，或者定义算子$T:H_0^1\to H_0^1$为：给定$h\in H_0^1$，令$Th\in H_0^2$是唯一满足 $\iu \lap(Th) \lap v=\iu \nab(Th)\cdot\nab v~(\forall v\in H_0^2)$的元素（由Lax--Milgram定理证明），验证它是自伴紧算子，于是得到一组正交基，再适当地选取 $u_0$ 在该正交基展开下的系数。
\end{problem}
\end{document}
